% TOP_PROBLEM: {"problemId": "problem.global-attractor-conjecture-complex-balanced-mass-action", "canonicalTitle": "Global Attractor Conjecture for Complex-Balanced Mass-Action Systems", "releaseRank": 297, "primaryDomain": "applied_computational_mathematics", "primaryDomainLabel": "Applied & computational mathematics", "displayStatus": "claimed_pending_refereeing", "releaseStatus": "claimed_pending_refereeing"}
% !TeX program = lualatex
% ATTEMPT_STATUS: unresolved
\documentclass[11pt]{article}
\usepackage[margin=25mm]{geometry}
\usepackage{fontspec}
\setmainfont{Latin Modern Roman}
\usepackage{amsmath,amssymb,mathtools}
\usepackage{unicode-math}
\setmathfont{Latin Modern Math}
\usepackage{xurl,hyperref}
\hypersetup{colorlinks=true,urlcolor=blue}
\setcounter{secnumdepth}{0}
\setlength{\parindent}{0pt}
\setlength{\parskip}{5pt}
\setlength{\emergencystretch}{3em}
\begin{document}
\section*{\texorpdfstring{Global Attractor Conjecture for Complex-Balanced Mass-Action Systems}{Global Attractor Conjecture for Complex-Balanced Mass-Action Systems}}\label{297}
\subsection{Definitions and mathematical statement}
A reaction $y\to y'$ has complexes $y,y'\in{\mathbb{Z}}_{\ge0}^d$ and rate $\kappa_{y\to y'}>0$. Mass action gives
\[
 \dot x=\sum_{y\to y'}\kappa_{y\to y'}x^y(y'-y),\qquad
 x^y=\prod_i x_i^{y_i}.
\]
Let $S$ be the real span of reaction vectors. A positive equilibrium $x^*$ is complex-balanced when, at each complex, the total incoming flux equals the outgoing flux at $x^*$. For every such system and $x_0>0$, the assertion is global existence and convergence to the unique positive equilibrium in
$(x_0+S)\cap{\mathbb{R}}_{>0}^d$. The entire compatibility class need not be bounded. The question is about all positive trajectories, not merely local stability near an equilibrium or convergence conditional on persistence.
\par\smallskip\textit{Formulation and concept references:} \href{https://arxiv.org/html/2501.06354v1}{[S1]}.

\subsection{Short English statement}
Must every positive trajectory of a complex-balanced mass-action system converge to the unique positive equilibrium in its stoichiometric class?
\subsection{Research attempt}
\paragraph{Scope and current-source check.}
This attempt was developed by GPT-6 Astra Ultra on 27 September 2026.
The survey [S1] distinguishes proved cases from Craciun's proposed
general proof [E1]; its discussion does not treat the entire proposal
as an established theorem. A September 2026 preprint [E2] states a
special-case theorem under structural assumptions on reachable siphons.
Those assumptions are not part of the catalogue's universal target.
Neither manuscript is imported here as a verified full solution.

The work proves global existence and boundedness, existence and
uniqueness of the positive equilibrium in each positive class, and
convergence under an explicit conservation-law condition on siphons.
It isolates persistence as the missing general step. The underlying
entropy mechanism is classical; the derivation and its boundary
limitations are made explicit rather than advertised as new.

\paragraph{The entropy and its exact dissipation inequality.}
Fix one complex-balanced positive equilibrium $x^*$, not necessarily
in the class of the initial point. For each reaction $y\to y'$ write
\[
 q_{y\to y'}=\kappa_{y\to y'}(x^*)^y>0,
 \qquad z_i=x_i/x_i^*,\qquad a_y=z^y.
\]
Complex balance says that the sum of the $q$'s entering any complex
equals their sum leaving it. Define
\[
 V(x)=\sum_{i=1}^d\left(x_i\log\frac{x_i}{x_i^*}-x_i+x_i^*\right),
                                                               \tag{1}
\]
with $0\log0=0$ on the closed orthant. Each summand is nonnegative
and tends to infinity as $x_i\to\infty$.
For positive $x$ along a solution,
\[
 \dot V=\sum_{y\to y'}q_{y\to y'}a_y\log(a_{y'}/a_y).
\]
The scalar inequality $a\log(b/a)\le b-a$, for $a,b>0$, follows
from $\log s\le s-1$, with equality exactly when $a=b$.
Consequently
\[
 \dot V\le\sum_{y\to y'}q_{y\to y'}(a_{y'}-a_y)=0.      \tag{2}
\]
The last equality is obtained by collecting the coefficient of each
$a_y$ and using complex balance. Moreover equality holds if and only if
\[
       a_y=a_{y'}\text{ on every reaction},
       \quad\text{equivalently}\quad \log(x/x^*)\in S^\perp.
                                                               \tag{3}
\]
The logarithm in (3) is componentwise. A condition on its sum alone
would lose the stoichiometric information.

\paragraph{Global existence despite unbounded compatibility classes.}
Since $V$ is continuous and coercive on the closed nonnegative orthant,
its sublevel set $\{V\le V(x_0)\}$ is compact. While the solution is
positive, (2) keeps it in this set, hence all coordinates are bounded
by a constant $M$ depending on the initial point and the system.

For a reaction which decreases coordinate $i$, its source has $y_i\ge1$,
so its monomial contains a factor $x_i$. On the bounded box,
all negative contributions to $\dot x_i$ are therefore bounded below
by $-C_i x_i$. Positive contributions can be discarded in this lower
bound, giving
\[
                       x_i(t)\ge x_i(0)e^{-C_it}>0        \tag{4}
\]
on any finite existence interval. A coordinate cannot reach zero in
finite time, and no coordinate can escape to infinity. The vector
field is polynomial and locally Lipschitz on a neighborhood of the
compact box, so the standard ODE continuation theorem gives existence
for every $t\ge0$.

Also $\dot x\in S$, so $x(t)\in x_0+S$.
This proves global positive existence and boundedness without assuming
the whole compatibility class is bounded. The lower bound (4) decays
as $t\to\infty$ and is not a uniform persistence estimate.

\paragraph{The positive equilibrium in each class.}
Consider the closed feasible set
$\overline P=(x_0+S)\cap\mathbb R_{\ge0}^d$.
Coercivity gives a minimizer $\bar x$ of $V$ on this set.
It cannot have a zero coordinate. Indeed the segment
$x(\epsilon)=(1-\epsilon)\bar x+\epsilon x_0$ lies in $\overline P$
and is positive for $\epsilon>0$. On coordinates where $\bar x_i=0$,
the change in (1) has leading term
$\epsilon x_{0,i}\log\epsilon$; on the other coordinates it is
$O(\epsilon)$. A nonempty set of zero coordinates would therefore give
$V(x(\epsilon))<V(\bar x)$ for small $\epsilon$, a contradiction.
Thus $\bar x>0$.

At this interior minimizer, differentiating along every vector in $S$
shows $\nabla V(\bar x)=\log(\bar x/x^*)\in S^\perp$.
Condition (3) then holds. In each connected component of the reaction
graph, the factors $a_y$ are constant, since they agree along every
edge. Multiplying the balanced fluxes $q$ in that component by the
common factor preserves complex balance. Hence $\bar x$ is itself
a complex-balanced equilibrium in the desired class.

For uniqueness, any positive equilibrium has $\dot V=0$, and hence
obeys (3). If $u,v$ were two in the same class, then
$u-v\in S$ and $\log u-\log v\in S^\perp$, implying
\[
                     \sum_i(u_i-v_i)(\log u_i-\log v_i)=0.
\]
Every summand is nonnegative and vanishes only when $u_i=v_i$,
because the logarithm is strictly increasing. Therefore $u=v$.
Denote this unique equilibrium by $\widehat x$.

\paragraph{Persistence would now give convergence.}
Suppose the positive trajectory satisfies
$\liminf_{t\to\infty}x_i(t)>0$ for every coordinate.
It is then eventually in a compact subset of the positive class.
The decreasing nonnegative function $V(x(t))$ has a limit.
Its omega-limit set is nonempty, compact and invariant, and $V$ is
constant on it by continuity. Along every trajectory in this set,
$\dot V=0$, so (3) holds and every point is the unique equilibrium
$\widehat x$. Thus the omega-limit set is a singleton and
\[
                                  x(t)\longrightarrow\widehat x.
                                                               \tag{5}
\]
These compactness and invariance statements follow from the polynomial
flow on a bounded set; no boundary differentiability of $V$ is used
because the set here is strictly positive. The missing hypothesis in
this paragraph is precisely persistence.

\paragraph{A concrete sufficient condition using siphons.}
A nonempty set $I$ of species is a siphon if every reaction producing
some species in $I$ has at least one species in $I$ in its source.
Assume that for every nonempty siphon $I$ there is a vector
\[
 w\in S^\perp\cap\mathbb R_{\ge0}^d,
 \qquad w\ne0,\qquad \operatorname{supp}(w)\subseteq I.  \tag{6}
\]
Then the trajectory is persistent and (5) follows. Here is a proof
that also explains why boundary equilibria alone are not the only
objects one should check in an a priori argument.

If a boundary point $z$ were in the omega-limit set, let
$I=\{i:z_i=0\}$, which is nonempty. The compact omega-limit set
contains a full trajectory through $z$, defined for positive and
negative time and staying in the nonnegative orthant. This can be
obtained by taking limits of the shifted trajectories $x(t_j+s)$
on each bounded $s$-interval and using uniqueness; boundedness of the
vector field gives the required equicontinuity.
Each zero coordinate has a local minimum at time zero on that full
trajectory, so its derivative is zero there.

At $z$, every reaction consuming such a coordinate has zero rate.
All remaining contributions to that coordinate's derivative are
nonnegative. Their sum being zero, each production rate must vanish.
Therefore every reaction producing a species of $I$ has a zero
coordinate in its source: $I$ is a siphon.
Choose $w$ from (6). Along the original trajectory,
$w\cdot x(t)=w\cdot x_0>0$, but $w\cdot z=0$ since its support is
contained in $I$. Continuity contradicts $z$ being a limit point.
There are thus no boundary omega-limit points. Compactness makes the
distance from the omega-limit set to every coordinate hyperplane
positive, which is exactly the needed persistence assertion.

This proves a full convergence theorem under the checkable extra
condition (6). It has not been proved that every complex-balanced
network satisfies that condition, and in fact it need not.

\paragraph{Two examples distinguish sufficiency from necessity.}
For the directed cycle
\[
              A\mathrel{\mathop{\longrightarrow}^{\alpha}}B
              \mathrel{\mathop{\longrightarrow}^{\beta}}C
              \mathrel{\mathop{\longrightarrow}^{\gamma}}A,
\]
the only nonempty siphon is the full species set. It supports the
conservation vector $(1,1,1)$. In the class $a+b+c=m>0$, the unique
positive equilibrium is
\[
 (a,b,c)=\left(\frac J\alpha,\frac J\beta,\frac J\gamma\right),
 \qquad J=\frac m{\alpha^{-1}+\beta^{-1}+\gamma^{-1}}.
\]
The three fluxes equal $J$, verifying complex balance directly.
The preceding theorem proves global convergence. The directed cycle
has no reverse reactions, so the argument does not assume detailed
balance.

For contrast, consider $2A\rightleftarrows3A$ with positive rates
$k,\ell$. Its equation is $\dot x=x^2(k-\ell x)$, with positive
complex-balanced equilibrium $x^*=k/\ell$ and stoichiometric space
$S=\mathbb R$. There is no nonzero conservation vector, so (6) fails
for the siphon $\{A\}$. Nevertheless every positive solution moves
monotonically toward $x^*$ and converges to it: below $x^*$ the
derivative is positive, above it negative, and uniqueness prevents
crossing the equilibrium. Thus (6) is sufficient, not necessary.

\paragraph{Why entropy alone does not exclude the boundary.}
The function (1) has finite value when coordinates vanish; a term
with $x_i=0$ contributes $x_i^*$, not infinity. Sublevel compactness
therefore bounds concentrations above but does not bound them below.
For the scalar example, $V(0)=x^*$ and sufficiently large initial
states have $V(x_0)>x^*$, so their entropy sublevel even contains the
boundary point. Its actual exclusion uses the vector field, not an
infinite entropy barrier. Applying the interior equality analysis
of (2) at a boundary point would also insert undefined logarithms.

\paragraph{Verification and final status.}
The local checks verify flux balance, the equilibrium formula, the
siphons of the directed cycle, and the conserved quantity. They test
the dissipation identity on rational sample states with numerical
logarithms only as a sign check. The existence, uniqueness, positivity,
and conditional convergence arguments above are analytic, not inferred
from time-stepping trajectories.

\textbf{UNRESOLVED.} Global bounded positive existence and the unique
positive class equilibrium are proved here, as is convergence under
(6). The first general gap is exclusion of boundary omega-limit points
without that conservation-law hypothesis. Specific next targets are a
boundary-repelling function for siphons lacking a supported conservation
vector, compatible estimates across intersecting boundary faces, and a
uniform persistence mechanism valid for all complex-balanced networks.
The recent source claims are not substituted for those missing proofs.

\subsection{Sources}
\begin{itemize}
\item[S1]Feliu–Shiu: From chemical reaction networks to algebraic and polyhedral geometry – and back again.
\url{https://arxiv.org/html/2501.06354v1}.
\textit{Formulation source.}
\item[E1]Gheorghe Craciun, \emph{Toric Differential Inclusions and a Proof of the Global Attractor Conjecture}, 2015 preprint.
\url{https://arxiv.org/abs/1501.02860}.
\textit{General proof proposal, whose status is qualified using [S1]; not an imported proof here.}
\item[E2]Carsten Wiuf, \emph{A Proof of the Global Attractor Conjecture in a Special Case}, submitted 21 September 2026.
\url{https://arxiv.org/abs/2609.24553}.
\textit{Recent restricted-siphon result; its additional hypotheses are not removed.}
\end{itemize}
\end{document}
